\documentclass{article}
\usepackage[UTF8]{ctex}
\usepackage{amsmath,amssymb,amsfonts}
\usepackage{enumitem}
\usepackage{geometry}
\geometry{a4paper, margin=1in}
% === 错题本排版优化（不参与题目标识） ===
\setlength{\parindent}{0pt}
\setlength{\parskip}{0.35em plus 0.10em minus 0.08em}
\setlength{\emergencystretch}{3em}
\setlist{itemsep=0.25em, topsep=0.40em, parsep=0pt}
\allowdisplaybreaks[2]
% === 错题本排版优化结束 ===
\begin{document}
\section*{上海交通大学2024年数学分析真题}
\section*{42．上海交通大学2024年数学分析试卷}

\subsection{题目1}
1．具有原函数的函数一定 Riemann 可积．

\subsection{题目2}
2．存在 \((-\infty,+\infty)\) 上没一点都不连续的函数．

\subsection{题目3}
3．若 \(\int_{a}^{+\infty} f(x) \mathrm{d} x\) 收敛，\(f(x)>0\) 且 \(f(x)\) 在 \([0,+\infty)\) 上连续，则 \(\lim _{x \rightarrow+\infty} f(x)=0\) ．

\subsection{题目4}
4．如果函数 \(f\) 在 \([a, b]\) 上可导，而导函数 \(f^{\prime}\) 在 \([a, b]\) 上不连续，那么 \(f^{\prime}\) 不一定能取到 \(f^{\prime}(a)\) 与 \(f^{\prime}(b)\)之间的一切值．

\subsection{题目5}
5．级数 \(\sum_{n=1}^{\infty} x^{\alpha} e^{-n x}(\alpha>1)\) 在 \([0,+\infty)\) 上一致收敛：

二．（25 分）计算题．

\begin{itemize}

\subsection{题目6}
  \item[1．] （12 分）设函数 \(f\) 在 \([1,+\infty)\) 上可导，如果 \(\lim _{x \rightarrow+\infty}\left[2 f(x)+x f^{\prime}(x)\right]=3\) ，求 \(\lim _{x \rightarrow+\infty} f(x)\) ．

\subsection{题目7}
  \item[2．] （13 分）计算 \(\int_{0}^{+\infty} e^{-x^{2}} \mathrm{~d} x\) ．
\end{itemize}

三。（25分）设 \(f(x)\) 为 \((-\infty,+\infty)\) 上的连续函数，周期函数，\(f(x)\) 不为常值函数．试证明 \(f(x)\) 存在最小正周期．

四．（20 分）求重积分

\[
I=\iint_{D} \sqrt{\frac{1-x^{2}-y^{2}}{1+x^{2}+y^{2}}} \mathrm{~d} x \mathrm{~d} y
\]

其中 \(D\) 是圆周 \(x^{2}+y^{2}=1\) 及坐标轴所围成的第一象限区域。

五．（25 分）设 \(u=u\left(x_{1}, x_{2}, \cdots, x_{n}\right)\) 是一个二阶连续可导的 \(n\) 元函数，\(u>0\) ，令

\[
\nabla=\left(\frac{\partial}{\partial x_{1}}, \frac{\partial}{\partial x_{2}}, \cdots, \frac{\partial}{\partial x_{n}}\right) .
\]

\(D^{2} u=\left(\frac{\partial^{2} u}{\partial x_{i} \partial x_{j}}\right)\) 为矩阵。同时

\[
|\nabla u|=\left(\sum_{i=1}^{n} u_{x_{i}}^{2}\right)^{\frac{1}{2}},\left|D^{2} u\right|=\left(\sum_{i, j=1}^{n} u_{x_{i} x_{j}}^{2}\right)^{\frac{1}{2}} .
\]

试证明：对任意的 \(\theta \in(0,1)\) 成立恒等式

\[
\left|D^{2} u\right|^{2}=\frac{1}{\theta^{2}} u^{2-2 \theta}\left|D^{2} u^{\theta}\right|^{2}-(1-\theta)^{2} u^{-2}|\nabla u|^{4}+2(1-\theta) u^{-1} \nabla u\cdot\left(D^{2}u\,\nabla u\right) .
\]

六．（25 分）设 \(n \geq 2\) 为整数，函数 \(f \in C^{2}([0, R])\) 满足 \(f_{r}(0)=0\) ．试证明

\[
\sup _{r \in(0, R)}\left|\frac{f_{r}(r)}{r}\right| \leq \frac{1}{n} \sup _{r \in(0, R)}\left|f_{r r}+\frac{n-1}{r} f_{r}\right| .
\]
\end{document}
